% =============================================================================
%  BÁO CÁO ĐỒ ÁN MÔN HỌC – LẬP TRÌNH CƠ SỞ DỮ LIỆU
%  Đề tài: HỆ THỐNG QUẢN LÝ PHÒNG KHÁM (QLPK)
%  Biên dịch:  xelatex BaoCao_QLPK.tex   (chạy 2 lần để cập nhật mục lục)
% =============================================================================
\documentclass[12pt,a4paper]{report}

% --------- ENGINE & NGÔN NGỮ -------------------------------------------------
\usepackage{fontspec}
\usepackage{polyglossia}
\setdefaultlanguage{vietnamese}
\setmainfont{DejaVu Serif}
\setsansfont{DejaVu Sans}
\setmonofont{DejaVu Sans Mono}[Scale=0.85]

% --------- TRÌNH BÀY ---------------------------------------------------------
\usepackage[a4paper,top=2.5cm,bottom=2.5cm,left=3cm,right=2cm]{geometry}
\usepackage{graphicx}
\usepackage{xcolor}
\usepackage{hyperref}
\usepackage{titlesec}
\usepackage{enumitem}
\usepackage{tabularx}
\usepackage{longtable}
\usepackage{array}
\usepackage{booktabs}
\usepackage{caption}
\usepackage{listings}
\usepackage{tcolorbox}
\usepackage{fancyhdr}
\usepackage{setspace}
\usepackage{tikz}
\usetikzlibrary{arrows.meta,positioning,shapes}
\usepackage{float}

\onehalfspacing

% Cấu hình hyperref
\hypersetup{
    colorlinks=true,
    linkcolor=black,
    urlcolor=blue,
    citecolor=blue,
    pdftitle={Báo cáo đồ án Quản lý Phòng khám},
    pdfauthor={Nhóm sinh viên},
    pdfsubject={Lập trình cơ sở dữ liệu}
}

% Header / footer
\pagestyle{fancy}
\fancyhf{}
\fancyhead[L]{\textit{Đồ án môn học – Lập trình cơ sở dữ liệu}}
\fancyhead[R]{\textit{Hệ thống QLPK}}
\fancyfoot[C]{\thepage}
\renewcommand{\headrulewidth}{0.4pt}

% Định dạng tiêu đề
\titleformat{\chapter}[hang]
  {\normalfont\Huge\bfseries\color{blue!60!black}}
  {\thechapter.}{0.5em}{}
\titleformat{\section}
  {\normalfont\Large\bfseries\color{blue!50!black}}
  {\thesection}{0.5em}{}
\titleformat{\subsection}
  {\normalfont\large\bfseries\color{blue!40!black}}
  {\thesubsection}{0.5em}{}

% Định dạng code SQL & C#
\definecolor{codebg}{RGB}{248,248,248}
\definecolor{codekw}{RGB}{0,0,180}
\definecolor{codestr}{RGB}{163,21,21}
\definecolor{codecmt}{RGB}{0,128,0}

\lstdefinelanguage{TSQL}{
  morekeywords={SELECT,FROM,WHERE,INSERT,UPDATE,DELETE,INTO,VALUES,
    CREATE,PROCEDURE,FUNCTION,TRIGGER,VIEW,TABLE,AS,BEGIN,END,IF,ELSE,
    DECLARE,SET,RETURN,RETURNS,JOIN,LEFT,RIGHT,INNER,OUTER,ON,GROUP,BY,
    ORDER,HAVING,UNION,WITH,GO,USE,ALTER,ADD,CONSTRAINT,FOREIGN,KEY,
    REFERENCES,PRIMARY,UNIQUE,CHECK,DEFAULT,NULL,NOT,AND,OR,IN,EXISTS,
    BIT,INT,FLOAT,DATE,DATETIME,NVARCHAR,VARCHAR,IDENTITY,TRANSACTION,
    COMMIT,ROLLBACK,TRY,CATCH,RAISERROR,SCOPE_IDENTITY,GETDATE,ISNULL,
    AFTER,INSTEAD,OF,INSERT,DELETE,SUM,COUNT,MONTH,YEAR,CAST,LIKE,
    STRING_AGG,SCOPE_IDENTITY,INSERTED,DELETED},
  sensitive=false,
  morecomment=[l]{--},
  morecomment=[s]{/*}{*/},
  morestring=[b]',
  morestring=[b]"
}

\lstset{
  backgroundcolor=\color{codebg},
  basicstyle=\ttfamily\small,
  keywordstyle=\color{codekw}\bfseries,
  stringstyle=\color{codestr},
  commentstyle=\color{codecmt}\itshape,
  numbers=left,
  numberstyle=\tiny\color{gray},
  stepnumber=1,
  numbersep=8pt,
  showspaces=false,
  showstringspaces=false,
  showtabs=false,
  frame=single,
  rulecolor=\color{gray!50},
  tabsize=2,
  captionpos=b,
  breaklines=true,
  breakatwhitespace=true,
  literate=
    {á}{{\'a}}1 {à}{{\`a}}1 {ả}{{\h{a}}}1 {ã}{{\~a}}1 {ạ}{{\d{a}}}1
    {â}{{\^a}}1 {ấ}{{\'{\^a}}}1 {ầ}{{\`{\^a}}}1 {ẩ}{{\h{\^a}}}1 {ẫ}{{\~{\^a}}}1 {ậ}{{\d{\^a}}}1
    {ă}{{\u{a}}}1 {ắ}{{\'{\u{a}}}}1 {ằ}{{\`{\u{a}}}}1 {ẳ}{{\h{\u{a}}}}1 {ẵ}{{\~{\u{a}}}}1 {ặ}{{\d{\u{a}}}}1
    {é}{{\'e}}1 {è}{{\`e}}1 {ẻ}{{\h{e}}}1 {ẽ}{{\~e}}1 {ẹ}{{\d{e}}}1
    {ê}{{\^e}}1 {ế}{{\'{\^e}}}1 {ề}{{\`{\^e}}}1 {ể}{{\h{\^e}}}1 {ễ}{{\~{\^e}}}1 {ệ}{{\d{\^e}}}1
    {í}{{\'i}}1 {ì}{{\`i}}1 {ỉ}{{\h{i}}}1 {ĩ}{{\~i}}1 {ị}{{\d{i}}}1
    {ó}{{\'o}}1 {ò}{{\`o}}1 {ỏ}{{\h{o}}}1 {õ}{{\~o}}1 {ọ}{{\d{o}}}1
    {ô}{{\^o}}1 {ố}{{\'{\^o}}}1 {ồ}{{\`{\^o}}}1 {ổ}{{\h{\^o}}}1 {ỗ}{{\~{\^o}}}1 {ộ}{{\d{\^o}}}1
    {ơ}{{\o{o}}}1 {ớ}{{\'{\o{o}}}}1 {ờ}{{\`{\o{o}}}}1 {ở}{{\h{\o{o}}}}1 {ỡ}{{\~{\o{o}}}}1 {ợ}{{\d{\o{o}}}}1
    {ú}{{\'u}}1 {ù}{{\`u}}1 {ủ}{{\h{u}}}1 {ũ}{{\~u}}1 {ụ}{{\d{u}}}1
    {ư}{{\o{u}}}1 {ứ}{{\'{\o{u}}}}1 {ừ}{{\`{\o{u}}}}1 {ử}{{\h{\o{u}}}}1 {ữ}{{\~{\o{u}}}}1 {ự}{{\d{\o{u}}}}1
    {ý}{{\'y}}1 {ỳ}{{\`y}}1 {ỷ}{{\h{y}}}1 {ỹ}{{\~y}}1 {ỵ}{{\d{y}}}1
    {đ}{{\dj}}1 {Đ}{{\DJ}}1
}

\lstdefinestyle{csharp}{
  language=[Sharp]C,
  morekeywords={var,async,await,partial,using,namespace,public,private,protected,
    internal,class,static,void,int,string,bool,float,decimal,double,DateTime,
    List,Dictionary,return,if,else,foreach,for,while,try,catch,finally,throw,new,this},
}

% =============================================================================
\begin{document}

% ------------------- TRANG BÌA -----------------------------------------------
\begin{titlepage}
\centering
\vspace*{0.5cm}
{\large\textbf{TRƯỜNG ĐẠI HỌC MỞ TP. HỒ CHÍ MINH}\par}
{\large\textbf{KHOA CÔNG NGHỆ THÔNG TIN}\par}
\vspace{1cm}
\includegraphics[width=0.25\textwidth]{logo.png}
\vspace{0.5cm}

\rule{\linewidth}{0.5mm}\\[0.4cm]
{\huge\bfseries BÁO CÁO ĐỒ ÁN MÔN HỌC\\[0.3cm]}
{\Large\bfseries LẬP TRÌNH CƠ SỞ DỮ LIỆU\par}
\rule{\linewidth}{0.5mm}\\[1cm]

{\Huge\bfseries\color{blue!60!black} ĐỀ TÀI:\\[0.3cm]
HỆ THỐNG QUẢN LÝ PHÒNG KHÁM\par}

\vfill

\begin{flushleft}
\begin{tabular}{ll}
\textbf{Giảng viên hướng dẫn:} & \dotfill \\[0.3cm]
\textbf{Lớp:} & \dotfill \\[0.3cm]
\textbf{Nhóm thực hiện:} & \dotfill \\[0.3cm]
\textbf{Thành viên:} & 1. \dotfill \\[0.2cm]
 & 2. \dotfill \\[0.2cm]
 & 3. \dotfill \\[0.2cm]
 & 4. \dotfill \\[0.3cm]
\end{tabular}
\end{flushleft}

\vfill
{\large TP. Hồ Chí Minh, tháng 05 năm 2026 \par}
\end{titlepage}

% ------------------- LỜI CẢM ƠN ----------------------------------------------
\thispagestyle{empty}
\chapter*{LỜI CẢM ƠN}
\addcontentsline{toc}{chapter}{Lời cảm ơn}

Trong suốt quá trình học tập và thực hiện đồ án môn \textit{Lập trình cơ sở dữ liệu}, nhóm chúng em đã nhận được rất nhiều sự hỗ trợ, hướng dẫn tận tình từ quý thầy cô và sự hợp tác hiệu quả giữa các thành viên trong nhóm. Nhóm xin gửi lời cảm ơn chân thành nhất đến giảng viên hướng dẫn đã cung cấp kiến thức nền tảng về T-SQL, ADO.NET, LINQ và Entity Framework, đồng thời định hướng kiến trúc đa lớp cho đồ án.

Do thời gian và kiến thức còn hạn chế, báo cáo không tránh khỏi những thiếu sót. Nhóm rất mong nhận được những góp ý quý báu của quý thầy cô để đồ án ngày càng hoàn thiện hơn.

\vspace{1cm}
\begin{flushright}
\textit{Nhóm thực hiện}
\end{flushright}

% ------------------- MỤC LỤC -------------------------------------------------
\tableofcontents
\listoffigures
\listoftables

% =============================================================================
\chapter{GIỚI THIỆU CHUNG}

\section{Tên đề tài}

\textbf{Hệ thống Quản lý Phòng khám (Quản Lý Phòng Khám -- QLPK)} là ứng dụng nghiệp vụ cho phép quản trị, lễ tân, bác sĩ, điều dưỡng và nhân viên thu ngân của một cơ sở khám chữa bệnh cấp nhỏ và vừa thực hiện toàn bộ các nghiệp vụ chuyên môn: tiếp nhận bệnh nhân, đặt và theo dõi lịch hẹn, khám bệnh, kê toa thuốc, lập hóa đơn thanh toán, lập báo cáo doanh thu và báo cáo sử dụng thuốc.

\section{Thông tin nhóm thực hiện}

\begin{table}[H]
\centering
\renewcommand{\arraystretch}{1.4}
\caption{Danh sách thành viên và phân công nhiệm vụ}
\begin{tabularx}{\textwidth}{|c|l|l|X|}
\hline
\textbf{STT} & \textbf{Họ và tên} & \textbf{MSSV} & \textbf{Nhiệm vụ} \\
\hline
1 & \dotfill & \dotfill & Phân tích yêu cầu, thiết kế CSDL, viết T-SQL Scripts (View, Procedure, Function, Trigger). \\
\hline
2 & \dotfill & \dotfill & Lập trình lớp DAL (ADO.NET), lớp DTO, quản lý kết nối CSDL. \\
\hline
3 & \dotfill & \dotfill & Lập trình lớp BUS, áp dụng LINQ to Objects để xử lý nghiệp vụ. \\
\hline
4 & \dotfill & \dotfill & Lập trình giao diện Windows Forms (Guna UI), tích hợp báo cáo PDF và gửi mail nhắc lịch. \\
\hline
\end{tabularx}
\end{table}

\section{Mục tiêu đồ án}

Đồ án nhằm vận dụng tổng hợp các kiến thức đã học, cụ thể:

\begin{itemize}[leftmargin=*]
\item Phân tích, thiết kế cơ sở dữ liệu quan hệ trên Microsoft SQL Server cho một bài toán nghiệp vụ thực tế.
\item Thành thạo \textbf{T-SQL}: viết các View, Stored Procedure, User-defined Function, Trigger và quản lý Transaction để bảo đảm toàn vẹn dữ liệu.
\item Xây dựng ứng dụng \textbf{.NET Framework 4.7.2} theo kiến trúc \textbf{3 lớp} (Presentation – Business – Data Access) kết hợp lớp đối tượng truyền dữ liệu (DTO).
\item Sử dụng \textbf{ADO.NET} (SqlConnection, SqlCommand, SqlDataReader, Parameter) để truy xuất CSDL một cách an toàn.
\item Áp dụng \textbf{LINQ to Objects} để xử lý các tập hợp dữ liệu trong tầng nghiệp vụ; tham khảo \textbf{LINQ to XML} cho việc xuất / nhập dữ liệu phụ trợ.
\item Khảo sát các công nghệ mở rộng: \textbf{Entity Framework} (Database First) và \textbf{ASP.NET Core Web API} làm hướng phát triển trong tương lai.
\end{itemize}

\section{Đường link GitHub}

\begin{tcolorbox}[colback=blue!5,colframe=blue!50!black,title=Mã nguồn dự án]
\textbf{Repository:} \url{https://github.com/<tên-nhóm>/QuanLyPhongKham}\\
\textbf{Nhánh chính:} \texttt{main}\\
\textbf{Hướng dẫn cài đặt:} xem tệp \texttt{README.md}.
\end{tcolorbox}

% =============================================================================
\chapter{PHÂN TÍCH VÀ THIẾT KẾ CƠ SỞ DỮ LIỆU}

\section{Khảo sát nghiệp vụ}

Phòng khám hoạt động theo quy trình: bệnh nhân được lễ tân tiếp nhận và đặt lịch hẹn → đến ngày hẹn, bác sĩ tạo \emph{Phiếu khám bệnh}, ghi triệu chứng và chẩn đoán → bác sĩ kê \emph{Toa thuốc} gồm nhiều \emph{Chi tiết đơn thuốc} (có số lượng và đơn vị) → nhân viên thu ngân lập \emph{Hóa đơn} căn cứ trên tiền thuốc (tính từ toa) và tiền khám (lấy từ \emph{Dịch vụ}). Hệ thống lưu lịch sử khám của bệnh nhân, sinh báo cáo doanh thu theo ngày/tháng, báo cáo sử dụng thuốc theo tháng.

\section{Sơ đồ thực thể mối quan hệ (ERD)}

\begin{figure}[H]
\centering
% Đặt file ảnh ERD.png hoặc ERD.pdf cùng thư mục với báo cáo
\fbox{\parbox[c][6cm][c]{0.9\textwidth}{\centering
\textit{\textbf{[Chèn ảnh ERD tại đây]}}\\[0.3cm]
File ảnh: \texttt{images/ERD.png}\\
(Sử dụng \texttt{\textbackslash includegraphics[width=\textbackslash textwidth]\{images/ERD.png\}})}}
\caption{Sơ đồ thực thể mối quan hệ (ERD) của hệ thống QLPK}
\label{fig:erd}
\end{figure}

CSDL gồm \textbf{15 bảng} được phân nhóm như sau:

\begin{itemize}[leftmargin=*]
\item \textbf{Nhóm tài khoản \& phân quyền:} \texttt{Roles}, \texttt{TaiKhoan}, \texttt{AuditLog}.
\item \textbf{Nhóm bệnh nhân \& khám chữa bệnh:} \texttt{BenhNhan}, \texttt{LichHen}, \texttt{PhieuKhamBenh}, \texttt{Benh}, \texttt{ChuanDoan}.
\item \textbf{Nhóm thuốc \& toa:} \texttt{Thuoc}, \texttt{DonVi}, \texttt{CachDung}, \texttt{ToaThuoc}, \texttt{ChiTietDonThuoc}.
\item \textbf{Nhóm dịch vụ \& hóa đơn:} \texttt{DichVu}, \texttt{HoaDon}.
\end{itemize}

\subsection{Mô tả chi tiết các bảng}

\begin{longtable}{|p{3.2cm}|p{3cm}|p{8cm}|}
\hline
\textbf{Bảng} & \textbf{Khóa chính} & \textbf{Mô tả} \\
\hline\endhead
\texttt{TaiKhoan} & \texttt{maTaiKhoan} & Tài khoản người dùng (Lễ tân, Bác sĩ, Điều dưỡng, Thu ngân, Quản trị). Có ràng buộc \texttt{UNIQUE} trên \texttt{userName}. \\\hline
\texttt{Roles} & \texttt{maRole} & Vai trò trong hệ thống. \\\hline
\texttt{BenhNhan} & \texttt{maBenhNhan} & Hồ sơ bệnh nhân. Có cờ \texttt{isDeleted} cho \emph{soft-delete}, ràng buộc \texttt{UNIQUE} trên \texttt{CCCD}. \\\hline
\texttt{LichHen} & \texttt{maLichHen} & Lịch hẹn khám. Liên kết tới bác sĩ, điều dưỡng, bệnh nhân; có trạng thái mặc định \emph{Chờ khám}. \\\hline
\texttt{PhieuKhamBenh} & \texttt{maPKB} & Phiếu khám gồm ngày khám, triệu chứng, ngày tái khám, cờ đã gửi mail nhắc. \\\hline
\texttt{Benh}, \texttt{ChuanDoan} & \texttt{maBenh} / khóa kép & Danh mục bệnh và chẩn đoán cụ thể trên phiếu khám. \\\hline
\texttt{Thuoc} & \texttt{maThuoc} & Danh mục thuốc với đơn giá, đơn vị, cách dùng, số lượng tồn kho. \\\hline
\texttt{DonVi}, \texttt{CachDung} & \texttt{maDonVi} / \texttt{maCachDung} & Danh mục đơn vị (viên, hộp, ống…) và cách dùng. \\\hline
\texttt{ToaThuoc} & \texttt{maToaThuoc} & Toa thuốc gắn với phiếu khám. \\\hline
\texttt{ChiTietDonThuoc} & (\texttt{maThuoc}, \texttt{maToaThuoc}) & Bảng kết nối nhiều–nhiều giữa toa và thuốc, có cột \texttt{soLuong}. \\\hline
\texttt{DichVu} & \texttt{maDichVu} & Bảng giá các dịch vụ (Khám bệnh, Tiểu phẫu…). \\\hline
\texttt{HoaDon} & \texttt{maHoaDon} & Hóa đơn thanh toán, gồm tiền thuốc, tiền khám, tổng tiền. \\\hline
\texttt{AuditLog} & \texttt{logID} & Nhật ký thao tác phục vụ kiểm toán. \\\hline
\caption{Mô tả các bảng trong CSDL \texttt{QLPK}}
\end{longtable}

\subsection{Các ràng buộc dữ liệu}

Ngoài khóa chính và khóa ngoại (17 quan hệ FK), CSDL cài đặt các \texttt{CHECK CONSTRAINT}:

\begin{itemize}[leftmargin=*]
\item \texttt{chk\_SoLuongKe}: số lượng thuốc kê trên mỗi đơn phải $> 0$.
\item \texttt{chk\_SoLuongThuoc}: tồn kho thuốc phải $\geq 0$.
\item \texttt{chk\_TienDichVu}: giá dịch vụ phải $\geq 0$.
\end{itemize}

Đồng thời, các giá trị mặc định (DEFAULT) được khai báo cho các cột thời gian (\texttt{ngayLapHoaDon}, \texttt{ngayKham}, \texttt{ngayKeToa}) và cờ trạng thái (\texttt{isDeleted}, \texttt{DaGuiMail}, \texttt{trangThai}).

% -----------------------------------------------------------------------------
\section{Lập trình phía máy chủ (T-SQL Scripts)}

\subsection{View truy vấn dữ liệu}

Toàn bộ thao tác đọc của ứng dụng đều đi qua \textbf{View} nhằm tách biệt cấu trúc bảng vật lý khỏi tầng ứng dụng, đồng thời mã hóa sẵn các phép JOIN phức tạp. Hệ thống cài đặt \textbf{13 View} chính:

\begin{longtable}{|l|p{8.5cm}|}
\hline
\textbf{View} & \textbf{Mục đích} \\\hline\endhead
\texttt{v\_DanhSachBenhNhan} & Liệt kê bệnh nhân chưa bị xóa mềm (\texttt{isDeleted = 0}). \\\hline
\texttt{v\_DanhSachTaiKhoan} & JOIN \texttt{TaiKhoan} với \texttt{Roles} để lấy tên vai trò. \\\hline
\texttt{v\_DanhSachLichHen} & JOIN lịch hẹn với bệnh nhân, bác sĩ và điều dưỡng. \\\hline
\texttt{v\_DanhSachThuoc} & JOIN thuốc với đơn vị, cách dùng. \\\hline
\texttt{v\_DanhSachDichVu}, \texttt{v\_DanhSachBenh}, \texttt{v\_DanhSachDonVi}, \texttt{v\_DanhSachCachDung}, \texttt{v\_DanhSachRoles}, \texttt{v\_DanhSachChuanDoan} & Liệt kê danh mục tham chiếu. \\\hline
\texttt{v\_ChiTietPhieuKham} & Chi tiết phiếu khám kèm tên bệnh nhân, tên bác sĩ. \\\hline
\texttt{v\_ChiTietDonThuoc} & Chi tiết các thuốc trong một toa, kèm ngày kê. \\\hline
\texttt{v\_DanhSachHoaDon} & JOIN hóa đơn với phiếu khám, bệnh nhân, thu ngân để hiển thị tổng hợp. \\\hline
\caption{Danh sách View trong CSDL}
\end{longtable}

Ví dụ \texttt{v\_DanhSachThuoc}:

\begin{lstlisting}[language=TSQL,caption={View v\_DanhSachThuoc – chuẩn hóa truy vấn danh mục thuốc}]
CREATE VIEW [dbo].[v_DanhSachThuoc] AS
SELECT  t.maThuoc, t.tenThuoc, t.donGia, t.soLuong,
        t.maDonVi, t.maCachDung,
        dv.tenDonVi, cd.tenCachDung
FROM    Thuoc t
JOIN    DonVi dv     ON t.maDonVi    = dv.maDonVi
JOIN    CachDung cd  ON t.maCachDung = cd.maCachDung;
\end{lstlisting}

\subsection{User-defined Function}

Hệ thống cài đặt 2 \textbf{Function} phục vụ tính toán nghiệp vụ:

\paragraph{(1) Hàm vô hướng \texttt{f\_TinhTienThuoc}} tính tổng tiền thuốc của một Phiếu khám:

\begin{lstlisting}[language=TSQL,caption={Scalar Function tính tiền thuốc}]
CREATE FUNCTION dbo.f_TinhTienThuoc(@maPKB INT)
RETURNS FLOAT
AS
BEGIN
    DECLARE @TongTien FLOAT;
    SELECT  @TongTien = SUM(th.donGia * ctdt.soLuong)
    FROM    ToaThuoc tt
    JOIN    ChiTietDonThuoc ctdt ON tt.maToaThuoc = ctdt.maToaThuoc
    JOIN    Thuoc th             ON ctdt.maThuoc  = th.maThuoc
    WHERE   tt.maPKB = @maPKB;
    RETURN ISNULL(@TongTien, 0);
END;
\end{lstlisting}

\paragraph{(2) Table-Valued Function \texttt{f\_LichSuKhamBenh}} trả về toàn bộ lịch sử khám của một bệnh nhân, kèm danh sách thuốc đã kê và tổng chi phí:

\begin{lstlisting}[language=TSQL,caption={Table-Valued Function tra cứu lịch sử khám}]
CREATE FUNCTION dbo.f_LichSuKhamBenh(@maBenhNhan INT)
RETURNS TABLE
AS RETURN
(
    SELECT pkb.maPKB, pkb.ngayKham, pkb.trieuChung,
           STRING_AGG(th.tenThuoc, ', ') AS ToaThuoc,
           (dbo.f_TinhTienThuoc(pkb.maPKB) + ISNULL(hd.tienKham,0)) AS TongChiPhi
    FROM   PhieuKhamBenh pkb
    LEFT JOIN ToaThuoc tt          ON pkb.maPKB     = tt.maPKB
    LEFT JOIN ChiTietDonThuoc ctdt ON tt.maToaThuoc = ctdt.maToaThuoc
    LEFT JOIN Thuoc th             ON ctdt.maThuoc  = th.maThuoc
    LEFT JOIN HoaDon hd            ON pkb.maPKB     = hd.maPKB
    WHERE  pkb.maBenhNhan = @maBenhNhan
    GROUP BY pkb.maPKB, pkb.ngayKham, pkb.trieuChung, hd.tienKham
);
\end{lstlisting}

\subsection{Stored Procedure}

Toàn bộ thao tác Insert/Update/Delete đều được đóng gói trong \textbf{Stored Procedure} để chống SQL Injection và tách rõ logic giữa tầng DAL và CSDL. Hệ thống cài đặt 25 procedure, chia theo các nhóm:

\begin{itemize}[leftmargin=*]
\item \textbf{Đăng nhập \& Tài khoản:} \texttt{sp\_KiemTraDangNhap}, \texttt{sp\_ThemTaiKhoan}, \texttt{sp\_SuaTaiKhoan}, \texttt{sp\_XoaTaiKhoan}.
\item \textbf{Quản lý bệnh nhân:} \texttt{sp\_ThemBenhNhan}, \texttt{sp\_SuaBenhNhan}, \texttt{sp\_XoaBenhNhan} (soft-delete), \texttt{sp\_TimKiemBenhNhan}.
\item \textbf{Quản lý bệnh:} \texttt{sp\_ThemBenh}, \texttt{sp\_SuaBenh}, \texttt{sp\_XoaBenh}.
\item \textbf{Quản lý thuốc:} \texttt{sp\_ThemThuoc}, \texttt{sp\_SuaThuoc}, \texttt{sp\_XoaThuoc}.
\item \textbf{Lịch hẹn:} \texttt{sp\_ThemLichHen}, \texttt{sp\_CapNhatTrangThaiLichHen}.
\item \textbf{Phiếu khám \& toa thuốc (có Transaction):} \texttt{sp\_TaoPhieuKhamVaToaThuoc}, \texttt{sp\_ThemChuanDoan}, \texttt{sp\_KeThuocVaoDon}, \texttt{sp\_XoaThuocKhoiDon}.
\item \textbf{Hóa đơn (có Transaction):} \texttt{sp\_LapHoaDonThanhToan}.
\item \textbf{Dịch vụ:} \texttt{sp\_ThemDichVu}, \texttt{sp\_SuaDichVu}, \texttt{sp\_XoaDichVu}.
\item \textbf{Báo cáo:} \texttt{sp\_BaoCaoSuDungThuoc}, \texttt{sp\_DemSoLanKeThuoc}.
\item \textbf{Khác:} \texttt{sp\_CapNhatDaGuiMail}, \texttt{sp\_ThemToaThuoc}.
\end{itemize}

\subsection{Quản lý Giao tác (Transaction)}

Hai nghiệp vụ \emph{Tạo phiếu khám + toa thuốc} và \emph{Lập hóa đơn} là nghiệp vụ \textbf{đa câu lệnh INSERT/UPDATE} nên được bao trong \texttt{BEGIN TRANSACTION ... COMMIT/ROLLBACK} kèm khối \texttt{TRY-CATCH}:

\begin{lstlisting}[language=TSQL,caption={Stored procedure sp\_LapHoaDonThanhToan – sử dụng Transaction}]
CREATE PROCEDURE dbo.sp_LapHoaDonThanhToan
    @maPKB INT,
    @maTaiKhoan INT
AS
BEGIN
    SET NOCOUNT ON;
    BEGIN TRANSACTION;
    BEGIN TRY
        DECLARE @tienThuoc FLOAT = dbo.f_TinhTienThuoc(@maPKB);

        DECLARE @tienKham FLOAT;
        SELECT @tienKham = tienDichVu FROM DichVu WHERE tenDichVu = N'Khám bệnh';
        IF @tienKham IS NULL SET @tienKham = 50000;

        DECLARE @tongTien   FLOAT = @tienThuoc + @tienKham;
        DECLARE @ngayTaiKham DATE;
        SELECT @ngayTaiKham = ngayTaiKham FROM PhieuKhamBenh WHERE maPKB = @maPKB;

        INSERT INTO HoaDon(ngayLapHoaDon, tienThuoc, tienKham, tongTien,
                           ngayTaiKham, maPKB, maTaiKhoan)
        VALUES (GETDATE(), @tienThuoc, @tienKham, @tongTien,
                @ngayTaiKham, @maPKB, @maTaiKhoan);

        COMMIT TRANSACTION;
    END TRY
    BEGIN CATCH
        ROLLBACK TRANSACTION;
        DECLARE @msg NVARCHAR(4000) = ERROR_MESSAGE();
        RAISERROR(@msg, 16, 1);
    END CATCH
END;
\end{lstlisting}

Ý nghĩa: việc tính tiền thuốc và tạo hóa đơn được coi là một \emph{đơn vị nguyên tử} – hoặc thành công toàn bộ, hoặc khôi phục về trạng thái cũ. Nếu Trigger \texttt{trg\_CapNhatSoLuongThuocTrongKho} báo lỗi kho âm, transaction sẽ \texttt{ROLLBACK}, không sinh ra hóa đơn lỗi.

\subsection{Trigger}

Hệ thống cài đặt 4 \textbf{Trigger} để đảm bảo ràng buộc dữ liệu tự động:

\begin{enumerate}[leftmargin=*]
\item \textbf{\texttt{trg\_CapNhatSoLuongThuocTrongKho}} – \texttt{AFTER INSERT} trên \texttt{ChiTietDonThuoc}: kiểm tra tồn kho và tự động trừ kho khi kê thuốc.

\item \textbf{\texttt{trg\_HoanTraSoLuongThuoc}} – \texttt{AFTER DELETE} trên \texttt{ChiTietDonThuoc}: hoàn trả số lượng tồn kho khi xóa thuốc khỏi đơn.

\item \textbf{\texttt{trg\_KiemTraNgayDatLich}} – \texttt{AFTER INSERT, UPDATE} trên \texttt{LichHen}: chặn việc đặt lịch trong quá khứ.

\item \textbf{\texttt{trg\_ChanXoaTaiKhoan}} – \texttt{INSTEAD OF DELETE} trên \texttt{TaiKhoan}: không cho xóa tài khoản đã có lịch sử khám/hóa đơn.
\end{enumerate}

\begin{lstlisting}[language=TSQL,caption={Trigger kiểm tra và cập nhật tồn kho thuốc}]
CREATE TRIGGER trg_CapNhatSoLuongThuocTrongKho
ON ChiTietDonThuoc
AFTER INSERT
AS
BEGIN
    IF EXISTS (
        SELECT 1 FROM inserted i
        JOIN Thuoc t ON i.maThuoc = t.maThuoc
        WHERE t.soLuong < i.soLuong)
    BEGIN
        RAISERROR(N'Không đủ số lượng thuốc trong kho!', 16, 1);
        ROLLBACK TRANSACTION;
        RETURN;
    END

    UPDATE t
    SET    t.soLuong = t.soLuong - i.soLuong
    FROM   Thuoc t
    JOIN   inserted i ON t.maThuoc = i.maThuoc;
END;
\end{lstlisting}

% =============================================================================
\chapter{KIẾN TRÚC PHẦN MỀM}

\section{Mô hình 3 lớp (3-layer Architecture)}

Ứng dụng được tổ chức thành \textbf{4 project} trong cùng một solution \texttt{QuanLyPhongKham.sln}:

\begin{itemize}[leftmargin=*]
\item \textbf{\texttt{GUI\_QLPK}} (Presentation Layer) – Windows Forms, sử dụng thư viện UI hiện đại \emph{Guna.UI2.WinForms}.
\item \textbf{\texttt{QLPKBUS}} (Business Logic Layer) – Trung gian giữa GUI và DAL, kiểm tra ràng buộc nghiệp vụ.
\item \textbf{\texttt{QLPKDAL}} (Data Access Layer) – Tương tác trực tiếp với SQL Server qua ADO.NET.
\item \textbf{\texttt{QLPKDTO}} (Data Transfer Object) – Chứa lớp đối tượng truyền dữ liệu giữa các lớp.
\end{itemize}

\begin{figure}[H]
\centering
\begin{tikzpicture}[
    layer/.style={rectangle,draw=blue!60!black,thick,fill=blue!10,
                  minimum width=11cm,minimum height=1.4cm,align=center,
                  font=\bfseries\large},
    arrow/.style={-Stealth,thick,blue!60!black}
]
\node[layer]                       (gui) at (0, 5.4) {GUI\_QLPK\\(Presentation Layer)};
\node[layer]                       (bus) at (0, 3.0) {QLPKBUS\\(Business Logic Layer)};
\node[layer]                       (dal) at (0, 0.6) {QLPKDAL\\(Data Access Layer)};
\node[draw=red!60!black,thick,fill=red!10,minimum width=11cm,minimum height=1.2cm,align=center,
      font=\bfseries] (db) at (0,-1.8) {SQL Server – CSDL QLPK\\(Views, Stored Procedures, Functions, Triggers)};

\draw[arrow] (gui) -- (bus);
\draw[arrow] (bus) -- (dal);
\draw[arrow] (dal) -- (db);

\node[draw=green!60!black,thick,fill=green!10,minimum width=2.5cm,minimum height=8cm,
      align=center,font=\bfseries] (dto) at (7,1.8) {QLPKDTO\\(DTO)};
\draw[arrow,dashed,green!60!black] (gui.east) -- (dto.north west);
\draw[arrow,dashed,green!60!black] (bus.east) -- (dto.west);
\draw[arrow,dashed,green!60!black] (dal.east) -- (dto.south west);
\end{tikzpicture}
\caption{Mô hình 3 lớp + DTO của hệ thống QLPK}
\label{fig:3layer}
\end{figure}

\subsection{Presentation Layer – GUI\_QLPK}

\begin{itemize}[leftmargin=*]
\item Cung cấp các Form: \texttt{Login}, \texttt{QLPMMain}, \texttt{Home}, \texttt{DangKyKham}, \texttt{DSLaySoKham}, \texttt{QuanLyBenhNhan}, \texttt{ThemPhieuKhamBenh}, \texttt{KeToa}, \texttt{LapHoaDon}, \texttt{DanhSachHoaDon}, \texttt{BaoCaoDoanhThu}, \texttt{BaoCaoSuDungThuoc}, \texttt{QuanLyTaiKhoan}, \texttt{QuanLyThuoc}, \texttt{QuanLyDichVu}, \texttt{QuanLyLoaiBenh}, \texttt{QuanLyNhacHen}, \texttt{TraCuuBenhNhan}\ldots
\item Sử dụng \textbf{Guna.UI2} để tạo giao diện hiện đại, sử dụng \textbf{iText7} để xuất hóa đơn PDF, \textbf{SMTP} để gửi mail nhắc tái khám.
\item Tham chiếu trực tiếp \texttt{QLPKBUS} và \texttt{QLPKDTO}, \textbf{không trực tiếp gọi} đến \texttt{QLPKDAL}.
\end{itemize}

\subsection{Business Logic Layer – QLPKBUS}

\begin{itemize}[leftmargin=*]
\item Mỗi entity có một class BUS tương ứng: \texttt{BenhNhanBUS}, \texttt{HoadonBUS}, \texttt{PhieukhambenhBUS}, \texttt{ThuocBUS}, \texttt{ToathuocBUS}, \texttt{lichHenBUS}, \texttt{taiKhoanBUS}, \texttt{DichvuBUS}, \texttt{BenhBUS}, \texttt{ChandoanBUS}, \texttt{cachDungBUS}, \texttt{donviBUS}\ldots
\item Vai trò: kiểm tra ràng buộc, áp dụng LINQ to Objects để lọc / sắp xếp dữ liệu, đóng vai trò \textit{facade} cho lớp GUI.
\end{itemize}

\begin{lstlisting}[style=csharp,caption={Lớp BUS đóng vai trò trung gian giữa GUI và DAL}]
namespace QLPKBUS
{
    public class BenhNhanBUS
    {
        private BenhNhanDAL bnDAL;
        public BenhNhanBUS()  { bnDAL = new BenhNhanDAL(); }

        public bool   them (BenhNhanDTO bn) => bnDAL.Them(bn);
        public bool   sua  (BenhNhanDTO bn, int oldId) => bnDAL.Sua(bn, oldId);
        public bool   xoa  (BenhNhanDTO bn) => bnDAL.Xoa(bn);

        public List<BenhNhanDTO> select() => bnDAL.select();
        public List<BenhNhanDTO> selectByKeyWord(string kw) => bnDAL.SelectByKeyWord(kw);
    }
}
\end{lstlisting}

\subsection{Data Access Layer – QLPKDAL}

\begin{itemize}[leftmargin=*]
\item Mỗi entity có một class DAL: \texttt{benhnhanDAL}, \texttt{HoadonDAL}, \texttt{PhieukhambenhDAL}, \texttt{ThuocDAL}, \texttt{ToathuocDAL}, \texttt{lichHenDAL}, \texttt{taiKhoanDAL}, \texttt{ChandoanDAL}, \texttt{DichvuDAL}, \texttt{BenhDAL}, \texttt{donViDAL}, \texttt{cachDungDAL}, \texttt{loaiTaiKhoanDAL}, \texttt{ChiTietToaThuocDAL}.
\item Sử dụng \textbf{ADO.NET}: đối tượng \texttt{SqlConnection}, \texttt{SqlCommand}, \texttt{SqlDataReader}, \texttt{SqlParameter}; gọi \textbf{Stored Procedure} hoặc \textbf{View} đã định nghĩa trên SQL Server.
\item Chuỗi kết nối lấy từ \texttt{App.config} thông qua \texttt{ConfigurationManager.AppSettings}.
\end{itemize}

\subsection{Đối tượng trao đổi dữ liệu (DTO) – QLPKDTO}

Project DTO chứa các lớp POCO không phụ thuộc thư viện ADO.NET / WinForms, dùng để truyền dữ liệu giữa GUI \(\leftrightarrow\) BUS \(\leftrightarrow\) DAL: \texttt{BenhNhanDTO}, \texttt{hoadonDTO}, \texttt{phieukhambenhDTO}, \texttt{toathuocDTO}, \texttt{ChiTietToaThuocDTO}, \texttt{thuocDTO}, \texttt{dichvuDTO}, \texttt{lichHenDTO}, \texttt{taiKhoanDTO}, \texttt{loaiTaiKhoanDTO}, \texttt{benhDTO}, \texttt{chandoanDTO}, \texttt{donViDTO}, \texttt{cachdungDTO}.

\begin{lstlisting}[style=csharp,caption={DTO – ví dụ với lớp BenhNhanDTO}]
namespace QLPKDTO
{
    public class BenhNhanDTO
    {
        public int      MaBN              { get; set; }
        public string   TenBN             { get; set; }
        public string   DiachiBN          { get; set; }
        public string   GtBN              { get; set; }
        public DateTime NgsinhBN          { get; set; }
        public string   CanCuocCongDan    { get; set; }
        public string   Email             { get; set; }
    }
}
\end{lstlisting}

% =============================================================================
\chapter{THỰC THI KỸ THUẬT VÀ CÔNG NGHỆ}

\section{Truy xuất dữ liệu}

\subsection{ADO.NET}

ADO.NET là công nghệ truy cập dữ liệu chủ đạo của ứng dụng. Mọi truy vấn dữ liệu được thực hiện qua bộ ba lớp \texttt{SqlConnection}, \texttt{SqlCommand}, \texttt{SqlDataReader} hoặc \texttt{SqlDataAdapter} với đặc tả \textbf{Parameterized Query} để phòng chống SQL Injection. Chuỗi kết nối được khai báo trong file \texttt{App.config}:

\begin{lstlisting}[caption={Cấu hình chuỗi kết nối trong App.config}]
<appSettings>
  <add key="ConnectionString"
       value="Data Source=localhost;Initial Catalog=QLPK;Integrated Security=True"/>
</appSettings>
\end{lstlisting}

Ví dụ thao tác đọc danh sách bệnh nhân thông qua View và Stored Procedure tìm kiếm:

\begin{lstlisting}[style=csharp,caption={DAL gọi View và Stored Procedure qua ADO.NET}]
public List<BenhNhanDTO> select()
{
    var ds = new List<BenhNhanDTO>();
    using (var con = new SqlConnection(ConnectionString))
    using (var cmd = new SqlCommand("SELECT * FROM v_DanhSachBenhNhan", con))
    {
        cmd.CommandType = CommandType.Text;
        con.Open();
        using (var rd = cmd.ExecuteReader())
        {
            while (rd.Read())
            {
                ds.Add(new BenhNhanDTO {
                    MaBN           = (int)rd["maBenhNhan"],
                    TenBN          = rd["tenBenhNhan"].ToString(),
                    GtBN           = rd["gioiTinh"].ToString(),
                    NgsinhBN       = (DateTime)rd["ngaySinh"],
                    DiachiBN       = rd["diaChi"].ToString(),
                    CanCuocCongDan = rd["CCCD"].ToString(),
                    Email          = rd["email"] == DBNull.Value ? "" : rd["email"].ToString()
                });
            }
        }
    }
    return ds;
}

public List<BenhNhanDTO> SelectByKeyWord(string kw)
{
    using (var con = new SqlConnection(ConnectionString))
    using (var cmd = new SqlCommand("sp_TimKiemBenhNhan", con) {
            CommandType = CommandType.StoredProcedure })
    {
        cmd.Parameters.AddWithValue("@tuKhoa", kw);   // chống SQL Injection
        // ... ExecuteReader & ánh xạ DTO
    }
}
\end{lstlisting}

Thao tác gọi \textbf{Scalar Function} \texttt{f\_TinhTienThuoc} từ C\#:

\begin{lstlisting}[style=csharp,caption={Gọi User-defined Function từ ADO.NET}]
public decimal tienthuoc(hoadonDTO hd, int maPKB)
{
    using (var con = new SqlConnection(ConnectionString))
    using (var cmd = new SqlCommand("SELECT dbo.f_TinhTienThuoc(@maPKB)", con))
    {
        cmd.Parameters.AddWithValue("@maPKB", maPKB);
        con.Open();
        var result = cmd.ExecuteScalar();
        return result == DBNull.Value ? 0m : Convert.ToDecimal(result);
    }
}
\end{lstlisting}

\subsection{LINQ to Objects}

LINQ to Objects được áp dụng tại tầng BUS và GUI để xử lý các danh sách \texttt{List<DTO>} mà DAL trả về – ví dụ: lọc theo trạng thái, sắp xếp theo ngày, gom nhóm doanh thu theo ngày, đếm số bệnh nhân theo bác sĩ\ldots

\begin{lstlisting}[style=csharp,caption={Ví dụ áp dụng LINQ to Objects để gom nhóm hóa đơn theo ngày}]
var doanhThuTheoNgay = lsHoaDon
    .Where(h => h.NgayLapHoaDon.Month == thang && h.NgayLapHoaDon.Year == nam)
    .GroupBy(h => h.NgayLapHoaDon.Date)
    .Select(g => new {
        Ngay      = g.Key,
        TongTien  = g.Sum(x => x.TongTien),
        SoHoaDon  = g.Count()
    })
    .OrderBy(x => x.Ngay)
    .ToList();
\end{lstlisting}

\subsection{LINQ to XML}

LINQ to XML được tham khảo phục vụ chức năng \textbf{nhập / xuất danh mục thuốc} dưới dạng XML, để trao đổi dữ liệu với phần mềm thứ ba:

\begin{lstlisting}[style=csharp,caption={Xuất danh sách thuốc ra XML bằng LINQ to XML}]
var doc = new XDocument(
    new XDeclaration("1.0","utf-8","yes"),
    new XElement("DanhSachThuoc",
        from t in lsThuoc
        select new XElement("Thuoc",
            new XAttribute("Ma", t.MaThuoc),
            new XElement("Ten",     t.TenThuoc),
            new XElement("DonGia",  t.DonGia),
            new XElement("DonVi",   t.TenDonVi),
            new XElement("CachDung",t.TenCachDung),
            new XElement("SoLuong", t.SoLuong))));
doc.Save("DanhMucThuoc.xml");
\end{lstlisting}

\subsection{Entity Framework (định hướng mở rộng)}

Để chuẩn bị cho việc nâng cấp ứng dụng, nhóm đã khảo sát \textbf{Entity Framework} theo mô hình \emph{Database First}. Quy trình áp dụng:

\begin{enumerate}[leftmargin=*]
\item Thêm gói \texttt{EntityFramework} qua NuGet.
\item Sử dụng \emph{ADO.NET Entity Data Model} sinh tự động \texttt{QLPKEntities} từ CSDL \texttt{QLPK}.
\item Thay thế từng lớp DAL bằng truy vấn LINQ to Entities:
\end{enumerate}

\begin{lstlisting}[style=csharp,caption={Mẫu thay thế DAL bằng Entity Framework}]
using (var db = new QLPKEntities())
{
    var ds = db.BenhNhans
              .Where(b => !b.isDeleted)
              .OrderBy(b => b.tenBenhNhan)
              .Select(b => new BenhNhanDTO {
                  MaBN  = b.maBenhNhan,
                  TenBN = b.tenBenhNhan,
                  // ...
              }).ToList();
}
\end{lstlisting}

\section{Xây dựng Web API (định hướng mở rộng)}

Bước tiếp theo của đồ án là tách lớp DAL/BUS thành dịch vụ \textbf{ASP.NET Core Web API} để có thể phục vụ thêm cho client di động/web. Bảng các endpoint dự kiến:

\begin{table}[H]
\centering
\renewcommand{\arraystretch}{1.3}
\caption{Bảng các endpoint REST dự kiến}
\begin{tabularx}{\textwidth}{|l|l|X|}
\hline
\textbf{Method} & \textbf{Endpoint} & \textbf{Mô tả} \\\hline
GET  & \texttt{/api/benhnhan} & Lấy danh sách bệnh nhân chưa bị xóa. \\\hline
GET  & \texttt{/api/benhnhan/\{id\}} & Lấy bệnh nhân theo mã. \\\hline
POST & \texttt{/api/benhnhan} & Thêm mới bệnh nhân. \\\hline
PUT  & \texttt{/api/benhnhan/\{id\}} & Cập nhật. \\\hline
DELETE & \texttt{/api/benhnhan/\{id\}} & Xóa mềm. \\\hline
POST & \texttt{/api/phieukham} & Tạo phiếu khám + toa thuốc (gọi SP có transaction). \\\hline
POST & \texttt{/api/hoadon/lap} & Lập hóa đơn thanh toán. \\\hline
GET  & \texttt{/api/baocao/doanhthu?thang=\&nam=} & Báo cáo doanh thu. \\\hline
\end{tabularx}
\end{table}

Mẫu Controller cơ bản:

\begin{lstlisting}[style=csharp,caption={Mẫu Controller ASP.NET Core trả JSON}]
[ApiController]
[Route("api/[controller]")]
public class BenhNhanController : ControllerBase
{
    private readonly BenhNhanBUS _bus = new BenhNhanBUS();

    [HttpGet]
    public ActionResult<IEnumerable<BenhNhanDTO>> Get() => _bus.select();

    [HttpPost]
    public IActionResult Post([FromBody] BenhNhanDTO dto)
        => _bus.them(dto) ? Ok() : BadRequest();
}
\end{lstlisting}

Định dạng phản hồi mặc định là \textbf{JSON}; có thể hỗ trợ \textbf{XML} bằng cách thêm
\texttt{services.AddControllers().AddXmlSerializerFormatters();}.

% =============================================================================
\chapter{KẾT QUẢ VÀ ĐÁNH GIÁ}

\section{Demo ứng dụng}

Phần này trình bày một số màn hình tiêu biểu của hệ thống QLPK. Mỗi hình ảnh đặt trong thư mục \texttt{images/} và được tham chiếu bằng \texttt{\textbackslash includegraphics}.

\subsection{Đăng nhập}
\begin{figure}[H]
\centering
\fbox{\parbox[c][4cm][c]{0.9\textwidth}{\centering\textit{[Chèn ảnh: images/login.png]}}}
\caption{Màn hình đăng nhập – kiểm tra qua \texttt{sp\_KiemTraDangNhap}.}
\end{figure}

\subsection{Quản lý bệnh nhân}
\begin{figure}[H]
\centering
\fbox{\parbox[c][4cm][c]{0.9\textwidth}{\centering\textit{[Chèn ảnh: images/quanlybenhnhan.png]}}}
\caption{Form quản lý bệnh nhân với chức năng thêm, sửa, xóa mềm và tìm kiếm.}
\end{figure}

\subsection{Tạo phiếu khám và kê toa thuốc}
\begin{figure}[H]
\centering
\fbox{\parbox[c][4cm][c]{0.9\textwidth}{\centering\textit{[Chèn ảnh: images/ketoa.png]}}}
\caption{Form kê toa thuốc – Trigger sẽ trừ kho tự động khi thêm chi tiết đơn.}
\end{figure}

\subsection{Lập hóa đơn và in PDF}
\begin{figure}[H]
\centering
\fbox{\parbox[c][4cm][c]{0.9\textwidth}{\centering\textit{[Chèn ảnh: images/hoadon.png]}}}
\caption{Form lập hóa đơn – chạy SP có Transaction, xuất PDF qua iText7.}
\end{figure}

\subsection{Báo cáo doanh thu \& sử dụng thuốc}
\begin{figure}[H]
\centering
\fbox{\parbox[c][4cm][c]{0.9\textwidth}{\centering\textit{[Chèn ảnh: images/baocao\_doanhthu.png]}}}
\caption{Báo cáo doanh thu theo tháng – sử dụng \texttt{sp\_BaoCaoSuDungThuoc} \& LINQ.}
\end{figure}

\section{Đánh giá chuẩn đầu ra}

\begin{table}[H]
\centering
\renewcommand{\arraystretch}{1.4}
\caption{Tự đánh giá mức độ đạt chuẩn đầu ra của môn học}
\begin{tabularx}{\textwidth}{|l|X|c|}
\hline
\textbf{Kỹ năng} & \textbf{Minh chứng trong đồ án} & \textbf{Mức đạt} \\
\hline
T-SQL nâng cao & 13 View, 25 Stored Procedure, 2 Function, 4 Trigger, 2 SP có Transaction. & Tốt \\\hline
Thiết kế CSDL & 15 bảng chuẩn 3NF, 17 ràng buộc khóa ngoại, 3 ràng buộc CHECK, soft-delete. & Tốt \\\hline
ADO.NET & 14 lớp DAL, dùng Parameter, sử dụng View + SP cho mọi thao tác. & Tốt \\\hline
Kiến trúc 3 lớp + DTO & 4 project độc lập, tách biệt rõ trách nhiệm. & Tốt \\\hline
LINQ to Objects & Lọc, sắp xếp, gom nhóm dữ liệu báo cáo trong BUS/GUI. & Khá \\\hline
LINQ to XML & Cài đặt mẫu xuất danh mục thuốc, sẵn sàng tích hợp. & Khá \\\hline
Entity Framework & Khảo sát mô hình Database First, đã có kế hoạch chuyển đổi. & Trung bình \\\hline
ASP.NET Core Web API & Thiết kế danh sách endpoint, sẵn sàng cho giai đoạn 2. & Trung bình \\\hline
Kỹ năng làm việc nhóm & Phân chia rõ ràng, sử dụng Git để cộng tác. & Tốt \\\hline
\end{tabularx}
\end{table}

\section{Khó khăn và hướng phát triển}

\subsection{Khó khăn}
\begin{itemize}[leftmargin=*]
\item Việc đồng bộ giữa nhiều thành viên trên cùng CSDL ban đầu gặp trục trặc; nhóm đã khắc phục bằng cách tạo script \texttt{script.sql} làm chuẩn và đẩy lên GitHub.
\item Một số nghiệp vụ liên quan đến nhiều bảng (kê thuốc – trừ kho, lập hóa đơn – tính tiền) yêu cầu Transaction và Trigger phức tạp, dễ gây lỗi nếu thiếu \texttt{ROLLBACK} trong khối \texttt{CATCH}.
\item Mật khẩu hiện đang lưu plaintext trong bảng \texttt{TaiKhoan}; cần băm (hash) trong giai đoạn nâng cấp bảo mật.
\end{itemize}

\subsection{Hướng phát triển}
\begin{itemize}[leftmargin=*]
\item Thay thế lớp DAL hiện tại bằng \textbf{Entity Framework Core} để giảm boilerplate.
\item Tách phần dịch vụ thành \textbf{ASP.NET Core Web API} kèm xác thực \emph{JWT} và bổ sung client web bằng React.
\item Bổ sung mã hóa mật khẩu (BCrypt) và áp dụng nguyên tắc Bảo mật theo tầng (Auth Filter, Role-based Access).
\item Tích hợp thêm \emph{thanh toán điện tử} (VNPay/MoMo) và \emph{nhắc lịch} qua Zalo/SMS thay cho email.
\item Mở rộng \texttt{AuditLog} với Trigger \texttt{AFTER INSERT/UPDATE/DELETE} để ghi lại mọi thao tác quan trọng.
\end{itemize}

% =============================================================================
\chapter*{TÀI LIỆU THAM KHẢO}
\addcontentsline{toc}{chapter}{Tài liệu tham khảo}

\begin{enumerate}[leftmargin=*]
\item Microsoft Learn -- \textit{Transact-SQL reference}. \url{https://learn.microsoft.com/en-us/sql/t-sql/}
\item Microsoft Learn -- \textit{ADO.NET overview}. \url{https://learn.microsoft.com/en-us/dotnet/framework/data/adonet/}
\item Microsoft Learn -- \textit{LINQ (Language-Integrated Query)}. \url{https://learn.microsoft.com/en-us/dotnet/csharp/linq/}
\item Microsoft Learn -- \textit{LINQ to XML}. \url{https://learn.microsoft.com/en-us/dotnet/standard/linq/linq-xml-overview}
\item Microsoft Learn -- \textit{Entity Framework Database First}. \url{https://learn.microsoft.com/en-us/ef/ef6/modeling/designer/workflows/database-first}
\item Microsoft Learn -- \textit{Create web APIs with ASP.NET Core}. \url{https://learn.microsoft.com/en-us/aspnet/core/web-api/}
\item Microsoft Learn -- \textit{SQL Server Triggers}. \url{https://learn.microsoft.com/en-us/sql/relational-databases/triggers/dml-triggers}
\item Microsoft Learn -- \textit{Stored Procedures (Database Engine)}. \url{https://learn.microsoft.com/en-us/sql/relational-databases/stored-procedures/stored-procedures-database-engine}
\item Syncfusion -- \textit{Free .NET eBooks}. \url{https://www.syncfusion.com/ebooks}
\item iText7 -- \textit{Knowledge base}. \url{https://kb.itextpdf.com}
\item Guna.UI2 -- \textit{Documentation}. \url{https://docs.gunaui.com/}
\item Giáo trình \textit{Lập trình cơ sở dữ liệu} -- Khoa CNTT, Trường ĐH Mở TP.HCM, 2025.
\end{enumerate}

\end{document}
